% META: {"id":"1NSI-TC-CO-046","chapitre":"1NSI-TYPES-CONSTRUITS","type_objet":"corrige","capacites":["1NSI-TYPES-CONSTRUITS-C5"],"mode_creation":"ex_nihilo","sources_inspiration":[],"status":"approved","exercice_ref":"1NSI-TC-EX-046","origin":"needs_review","status_history":["needs_review","audited","accepted","approved"],"release_acceptance":"RELEASE_OWNER_BATCH_ACCEPTANCE","acceptance_closure_digest":"sha256:9b3ccf9a81c5520fbb7e03b7b2d3e7bf2057a02834b6d908bf3be5f49f3b553f","acceptance_receipt_digest":"sha256:4fad86f0282e0fa4e0419cca1ad6379ae7af5a280c04a4a90880f90a1c044242"}
% BEGIN-VERIFY
% a = [1, 2, 3]
% b = a
% b.append(4)
% assert a == [1, 2, 3, 4]
% assert a is b
% END-VERIFY
\begin{corrige}{1NSI-TC-EX-046}
\begin{enumerate}
  \item Le programme affiche :
\begin{console}
[1, 2, 3, 4]
True
\end{console}

  \item L'affectation \lstinline{b = a} ne crée pas une nouvelle liste : elle fait pointer \lstinline{b} vers \emph{le même objet} que \lstinline{a}. On dit que \lstinline{b} est un \textbf{alias} de \lstinline{a}. Toute modification via \lstinline{b} (ici \lstinline{b.append(4)}) modifie l'unique liste partagée, donc \lstinline{a} voit aussi le changement.

  \item L'opérateur \lstinline{is} teste si deux noms désignent \textbf{le même objet en mémoire} (même adresse), et non simplement si leurs valeurs sont égales. Ici \lstinline{a is b} vaut \lstinline{True} car \lstinline{a} et \lstinline{b} sont deux références vers la même liste.
\end{enumerate}
\end{corrige}
